%% ch09-troubleshooting.tex — JA4proxy Operator's User Guide
%% Chapter 9: Troubleshooting

\chapter{Troubleshooting}
\label{ch:troubleshooting}
\index{troubleshooting}

\newthought{Most \ja\ problems fall into one of five categories:}
containers that will not start, traffic not being blocked when it should be,
traffic being blocked when it should not be, performance problems, and
connectivity issues between components. This chapter covers all five.

Start with \code{docker compose ps} and \code{docker compose logs} for any
problem — most issues produce clear error messages in the container logs.

\section{Quick Diagnostic Commands}
\index{diagnostics}

These commands give you a fast overview of what is happening:

\begin{bashcode}
# Container status and health
docker compose ps

# All container logs, last 50 lines each
docker compose logs --tail=50

# Proxy-specific logs, live
docker compose logs -f ja4proxy

# Redis connectivity test
docker compose exec redis redis-cli PING

# Proxy health endpoint
curl -s http://localhost:8080/health | python3 -m json.tool

# HAProxy stats
curl -s http://localhost:8404/stats

# Active ban count
docker compose exec redis redis-cli --scan --pattern "ban:*" | wc -l

# Current dial setting
docker compose exec redis redis-cli GET "config:dial"
\end{bashcode}

\section{Symptom Lookup Table}
\index{troubleshooting!symptom table}

\subsection{Container Will Not Start}

\begin{table}[ht]
  \small
  \begin{tabularx}{\linewidth}{lXX}
    \toprule
    \textbf{Symptom} & \textbf{Likely cause} & \textbf{Fix} \\
    \midrule
    \code{Error: port 443 in use} & Another process is on port 443 & Stop the conflicting
      service or change the port mapping in \code{docker-compose.yml} \\
    \midrule
    \code{redis: connection refused} & Redis did not start before the proxy & Wait 10s
      and retry; Redis has a health check dependency. Run \code{docker compose restart ja4proxy} \\
    \midrule
    \code{No such file: proxy.yml} & Config file missing from expected path & Ensure
      \code{config/proxy.yml} exists. It is not auto-created; restore from \code{config/proxy.yml.example} \\
    \midrule
    \code{YAML parse error} & Syntax error in \code{config/proxy.yml} & Validate with
      \code{ja4p config validate} \\
    \midrule
    Build fails fetching Go modules & Network issue or proxy.golang.org unavailable & Try
      \code{docker compose build --no-cache proxy} with a working network connection \\
    \bottomrule
  \end{tabularx}
  \caption{Container startup failures}
\end{table}

\subsection{Proxy Running but Not Blocking}

\begin{table}[ht]
  \small
  \begin{tabularx}{\linewidth}{lXX}
    \toprule
    \textbf{Symptom} & \textbf{Likely cause} & \textbf{Fix} \\
    \midrule
    All traffic shows ALLOWED & Dial is 0 (monitor mode) & This is correct behaviour.
      Raise the dial in \code{config/proxy.yml} when ready \\
    \midrule
    Known-bad fingerprint not blocked & JA4 blacklist bypass disabled in config & Check
      \configkey{security\_policy.ja4\_blacklist\_bypass.enabled} \\
    \midrule
    Blacklisted fingerprint not blocked & Fingerprint not in Redis SET & After config
      reload, verify: \code{redis-cli SISMEMBER ja4:blacklist "fingerprint"} \\
    \midrule
    Country block not working & GeoIP database missing or stale & Run
      \code{make update-geoip} and reload config \\
    \midrule
    High-score traffic not blocked & Dial too low for score & Check the score in
      logs and compare to the block threshold at your current dial setting \\
    \bottomrule
  \end{tabularx}
  \caption{Traffic not being blocked as expected}
\end{table}

\subsection{Legitimate Traffic Being Blocked}

\begin{table}[ht]
  \small
  \begin{tabularx}{\linewidth}{lXX}
    \toprule
    \textbf{Symptom} & \textbf{Likely cause} & \textbf{Fix} \\
    \midrule
    All browser traffic blocked & ALPN bypass disabled & Set
      \configkey{security\_policy.alpn\_browser\_bypass.enabled: true} and reload \\
    \midrule
    Specific legitimate tool blocked & Tool fingerprint blacklisted & Move fingerprint
      from blacklist to whitelist, or remove from blacklist \\
    \midrule
    API client rate-limited & Rate limit too aggressive for API use case & Add
      client fingerprint to whitelist or increase rate limit thresholds \\
    \midrule
    Customer reports blocked & IP in wrong GeoIP country & Check the country
      shown in the log entry; update GeoIP if stale \\
    \midrule
    Monitoring tool blocked & Internal IP not in allowlist & Add monitoring agent's
      IP to \configkey{security.static\_ip\_allowlist} \\
    \bottomrule
  \end{tabularx}
  \caption{False positives — legitimate traffic blocked}
\end{table}

\subsection{High Latency}

\begin{table}[ht]
  \small
  \begin{tabularx}{\linewidth}{lXX}
    \toprule
    \textbf{Symptom} & \textbf{Likely cause} & \textbf{Fix} \\
    \midrule
    P95 connection latency $>$5ms & Redis TCP round-trips & Switch to Unix domain
      socket: set \configkey{redis.unix\_socket} in config \\
    \midrule
    P99 latency spikes & AbuseIPDB API timeout & AbuseIPDB calls are fire-and-forget;
      they should not affect latency. Check if the AbuseIPDB timeout is set correctly \\
    \midrule
    High latency on first connection from new IP & DNS lookup blocking & DNS enrichment
      uses async non-blocking lookups; verify \code{dns\_enrichment.async: true} \\
    \midrule
    Redis latency $>$2ms & Redis container resource-starved & Check
      \code{docker stats} for Redis CPU and memory; increase Docker memory limit \\
    \bottomrule
  \end{tabularx}
  \caption{Latency problems}
\end{table}

\subsection{Redis Connection Errors}

\begin{table}[ht]
  \small
  \begin{tabularx}{\linewidth}{lXX}
    \toprule
    \textbf{Symptom} & \textbf{Likely cause} & \textbf{Fix} \\
    \midrule
    Proxy logs: \code{Redis connection lost} & Redis container restarted & Proxy
      reconnects automatically; verify Redis is healthy \\
    \midrule
    \code{READONLY error} & Redis has entered read-only mode & Check Redis logs;
      may indicate disk-full condition \\
    \midrule
    \code{OOM command not allowed} & Redis hit \code{maxmemory} limit & Increase
      \code{maxmemory} in \code{config/redis.conf} or clear expired keys \\
    \midrule
    All metrics show zero bans & Redis data was flushed & Check if something
      ran \code{FLUSHALL}; this is the expected state after \code{make stop-clean} \\
    \bottomrule
  \end{tabularx}
  \caption{Redis connectivity and data issues}
\end{table}

\subsection{Metrics Not Appearing in Grafana}

\begin{table}[ht]
  \small
  \begin{tabularx}{\linewidth}{lXX}
    \toprule
    \textbf{Symptom} & \textbf{Likely cause} & \textbf{Fix} \\
    \midrule
    Grafana shows ``No data'' & Prometheus target is down & Check
      \url{http://localhost:9091/targets} — the ja4proxy target should be UP \\
    \midrule
    Prometheus target UP but no data & No traffic has been processed & Generate
      test traffic with \code{generate-tls-traffic.sh} \\
    \midrule
    Dashboard panels show different time ranges & Panel time range overrides set & Click
      each panel, edit, and remove time range overrides \\
    \midrule
    Loki shows no logs & Promtail not running & Check \code{docker compose ps}
      for the promtail service \\
    \bottomrule
  \end{tabularx}
  \caption{Observability issues}
\end{table}

\section{Log Analysis for Common Issues}
\index{log analysis}

\subsection{Finding the Root Cause of a Block}

When a connection is blocked and you want to know why:

\begin{bashcode}
# Search for a specific IP in logs
docker compose logs ja4proxy | grep "192.0.2.1"
\end{bashcode}

The \code{Reason:} field in the log line identifies the blocking cause:

\begin{table}[ht]
  \small
  \begin{tabularx}{\linewidth}{lX}
    \toprule
    \textbf{Reason string} & \textbf{Cause} \\
    \midrule
    \code{JA4 blacklisted} & Fingerprint is in the blacklist \\
    \code{Banned for 300s} & IP was auto-banned by rate limiter \\
    \code{Country blocked: RU} & Country is in the blacklist \\
    \code{TLS version rejected} & TLS 1.0 or 1.1 \\
    \code{Spamhaus DROP match} & IP is in Spamhaus DROP/EDROP list \\
    \code{Score: 85 exceeds threshold} & Risk score above block threshold at current dial \\
    \code{Rate limited by IP} & by\_ip strategy triggered \\
    \code{Rate limited by JA4} & by\_ja4 strategy triggered \\
    \bottomrule
  \end{tabularx}
  \caption{Block reason strings and their causes}
\end{table}

\subsection{Checking Signal Contributions}

At \code{DEBUG} log level, the proxy logs individual signal contributions. Enable
debug logging temporarily:

\begin{yamlcode}
logging:
  level: DEBUG   # Change from INFO to DEBUG
\end{yamlcode}

Reload config. Debug output includes each signal's score contribution:

\begin{lstlisting}[language=bash, numbers=none]
DEBUG | scorer | ip=185.220.0.1 | signal=abuseipdb | contrib=+30 | reason=score=75
DEBUG | scorer | ip=185.220.0.1 | signal=asn | contrib=+20 | reason=datacenter_asn
DEBUG | scorer | ip=185.220.0.1 | signal=beaconing | contrib=+40 | reason=cv=0.12
DEBUG | scorer | ip=185.220.0.1 | total_score=90
\end{lstlisting}

\begin{warningbox}[Debug logging volume]
At DEBUG level, every connection produces many log lines. On a busy proxy this
generates significant log volume. Enable debug logging only for targeted
investigations and disable it immediately after.
\end{warningbox}

\section{Getting Help}
\index{support}
\index{GitHub issues}

\begin{description}[leftmargin=3.5cm]
  \item[GitHub Issues] \url{https://github.com/seanpor/JA4proxy/issues} — for bug
    reports and feature requests. Include the output of \code{docker compose logs}
    and the relevant section of \code{config/proxy.yml} (with any credentials redacted).
  \item[JA4 Database] \url{https://ja4db.com} — for identifying unknown fingerprints.
  \item[JA4 Specification] \url{https://github.com/FoxIO-LLC/ja4} — for understanding
    how fingerprints are computed.
  \item[Reference Manual] The companion document for architecture detail, signal
    algorithm specifications, and API documentation.
\end{description}

When reporting an issue, include:
\begin{itemize}
  \item The symptom (what you expected vs what happened)
  \item The relevant section of \code{config/proxy.yml}
  \item Output of \code{docker compose ps} and \code{docker compose logs --tail=100}
  \item The \ja\ version (\code{git rev-parse --short HEAD})
  \item The host OS and Docker version
\end{itemize}
